\documentclass{article}
\usepackage[T2A]{fontenc}
\usepackage[utf8]{inputenc}
\usepackage{amsthm}
\usepackage{amsmath}
\usepackage{amssymb}
\usepackage{amsfonts}
\usepackage{mathrsfs}
\usepackage{fancyvrb}
\usepackage{indentfirst}
\usepackage[
  left=2cm, right=2cm, top=2cm, bottom=2cm, headsep=0.2cm, footskip=0.6cm, bindingoffset=0cm
]{geometry}
\usepackage[english,russian]{babel}


\begin{document}
\section*{Вариант 1}
В левой полости давление постоянно $p_0$, а в правой полости оно имеет постоянную составляющую и гармонически изменяющуюся по времени составляющую $p_0+p_1^+(\omega t)$. Закон изменения давления на правом торце представим в виде
\begin{equation}
    p_1^+=p_{1m}^+f_p(\omega t),    f_p(wt)=\exp(i \omega t),
\end{equation}
где $p_{1m}^+$ "---амплитуда пульсаций давления на торце канала; $\omega$ "--- частота пульсаций; $f_p(\omega t)$ "--- закон изменения давления.
Введем в рассмотрение безразмерные переменные
\begin{equation}
\begin{split}
        \psi = \delta_0/\ell \ll 1, \lambda = w_m/\delta_0 \ll 1, \Re = \delta_0^2\omega/v, \tau = \omega t, \xi = x/\ell, \zeta = z/\delta_0, \\
    V_z = w_m \omega U_\zeta, w = w_mW, u = u_mU, V_x = w_m \omega U_\xi \psi, \\
    p = p_0 + w_m\rho v \omega(\delta_0 \psi^2)^{-1}P, p^+ = w_m \rho v \omega(\delta_0 \psi^2)^{-1}P^+.
\end{split}
\end{equation}
Здесь $p$ "--- давление; $\rho, v$ "--- плотность и кинематографичесикй коэффициент вязкости жидкости; $V_x, V_z$ "--- проекции скорости движения жидкости на оси координат, $w_m$ "--- амплитуда прогиба пластины; $W$ "--- безразмерный прогиб пластины; $u_m$ "--- амплитуда продольного перемещения пластины; $U$ "--- безразмерное продольное перемещение пластины; $\psi, \lambda, \Re$ "--- параметры, характеризующие задачу.

\end{document}
